\section{Introduction}

La contaminación del aire se reconoce como un problema relevante de salud pública y, en contextos urbanos, una parte muy amplia de la población continúa expuesta a concentraciones que no cumplen con estándares orientados a la protección sanitaria \cite{PROB01}. En ese escenario, la comunicación del riesgo no puede descansar únicamente en series de concentraciones por contaminante, porque la lectura directa de esos datos resulta poco accesible para la población general y dificulta su uso por responsables públicos y operadores técnicos; por ello, los sistemas de índice de calidad del aire se consolidaron como mecanismos para traducir el estado de la atmósfera en mensajes comprensibles sobre riesgo y protección de la salud \cite{PROB02,PROB03}. Sin embargo, esa capa de interpretación sigue lejos de ser uniforme: ya desde las revisiones tempranas de índices de contaminación se documentaba una diversidad estructural considerable entre sistemas de cálculo y categorías reportadas \cite{PROB02}, y trabajos posteriores mostraron que los enfoques basados en un único contaminante dominante pueden subestimar la severidad del riesgo cuando varios contaminantes elevados coexisten o cuando el interés se desplaza hacia efectos agregados sobre salud \cite{PROB04,PROB05}. Esta limitación adquiere mayor peso cuando el monitoreo debe sostener decisiones urbanas más finas, en tiempo casi real y con mayor resolución espacial, porque los sistemas convencionales de vigilancia no escalan bien a esa necesidad y los despliegues basados en sensores de bajo costo, aunque prometedores, todavía arrastran problemas de cobertura, flexibilidad, calibración y confiabilidad metrológica \cite{PROB06,PROB07}.

En este marco, la lógica difusa aparece como una alternativa pertinente cuando el fenómeno a modelar no se presta fácilmente a descripciones estrictamente deterministas, ya que su formulación se apoya en variables lingüísticas, reglas condicionales difusas y mecanismos de inferencia aproximada orientados precisamente al tratamiento de sistemas complejos o mal definidos \cite{ART03}. Esa base conceptual ya se ha trasladado al dominio de calidad del aire en estudios recientes, aunque con propósitos y alcances distintos: se ha utilizado para clasificación difusa de calidad del aire interior \cite{ART04}, para control y automatización de ventilación en ambientes cerrados \cite{ART05}, para monitoreo IoT con clasificación en categorías ISPU/AQI \cite{ART06}, para evaluación en tiempo real de múltiples contaminantes en entornos de smart city \cite{ART07}, y para optimización de rutas y operación eficiente de arquitecturas WSN de monitoreo urbano \cite{ART08}. Precisamente esa amplitud de aplicaciones muestra que el uso de inferencia difusa en monitoreo de calidad del aire ya no es marginal, pero también evidencia que el campo se ha desarrollado con decisiones heterogéneas respecto a variables de entrada, tipo de sistema difuso, escala de despliegue, arquitectura tecnológica y forma de validación \cite{ART04,ART05,ART06,ART07,ART08}. Desde una perspectiva de síntesis, esa dispersión dificulta identificar qué configuraciones se repiten, qué decisiones de diseño son comparables y qué tan transferibles son los resultados reportados entre contextos de implementación \cite{ART04,ART05,ART06,ART07,ART08}.

Esta heterogeneidad justifica una revisión sistemática específica porque, en calidad del aire, el monitoreo adquiere valor cuando sus mediciones pueden apoyar la gestión pública, la protección sanitaria y la respuesta preventiva de la población \cite{ART01,ART02,PROB03}. Una síntesis rigurosa puede, por ello, aportar criterios más claros para la comparación académica de enfoques, ofrecer insumos a reguladores y planificadores que utilizan información atmosférica para orientar acciones de gestión ambiental \cite{ART01,PROB06}, y servir de apoyo al diseño de plataformas urbanas que demandan mayor resolución, flexibilidad operativa y capacidad de respuesta frente a variaciones locales del fenómeno \cite{PROB06,ART07,ART08}. La necesidad del estudio se refuerza porque la literatura reciente ya evidencia un campo activo y técnicamente diverso, pero todavía disperso en la forma de reportar variables, decisiones de modelado, arquitecturas e indicadores de desempeño \cite{ART04,ART05,ART06,ART07,ART08}. Desde la perspectiva metodológica, la revisión es viable porque el dominio cuenta con un conjunto identificable de estudios recientes, recuperables en bases indexadas y suficientemente homogéneos en tema como para permitir una síntesis estructurada siguiendo Kitchenham \cite{kitchenham2007}.

La literatura más próxima al tema ya ofrece dos líneas de antecedente que ayudan a situar esta revisión. La primera está formada por trabajos sobre AQI, desde revisiones amplias como las de Priti y Kumar \cite{ART01} y Suman \cite{ART02} hasta desarrollos aplicados de índice en contexto urbano como el de Murena \cite{PROB08}; entre ellos existe coincidencia en subrayar la necesidad de convertir concentraciones de contaminantes en información interpretable para la gestión y la protección de la salud, y en mostrar que esa traducción sigue afectada por diferencias en contaminantes considerados, métodos de cálculo, modos de agregación y categorías de interpretación. La segunda línea corresponde a estudios aplicados donde la lógica difusa se incorpora directamente al monitoreo y la evaluación de la calidad del aire. Allí aparecen puntos comunes claros, como el uso de reglas difusas para tratar incertidumbre y apoyar decisiones sobre estados de calidad del aire \cite{ART04,ART05,ART06,ART07}, pero también divergencias importantes en el tipo de entorno analizado, las variables de entrada, el objetivo del sistema y la arquitectura tecnológica adoptada: algunos trabajos se centran en clasificación \cite{ART04}, otros en control ambiental \cite{ART05}, otros en monitoreo IoT e interpretación de índices \cite{ART06,ART07}, y otros en la eficiencia operativa de redes WSN para ciudades inteligentes \cite{ART08}. La Tabla~\ref{tab:antecedentes_intro} resume estos antecedentes en función de su aporte principal, su cercanía con el problema estudiado y el aspecto que todavía dejan abierto frente a la presente revisión.

\begin{adjustwidth}{-\extralength}{0cm}
\begin{longtblr}[
  label={tab:antecedentes_intro},
  caption={Antecedentes que delimitan el alcance de la revisión sistemática.}
]{
  colspec={X[0.8,c,m] X[1.8,l,m] X[2.4,l,m] X[2.5,l,m]},
  width =\fulllength,
  rowhead=1,
  row{1}={font=\bfseries}
}
\toprule
Referencia & Tipo/aporte principal & Relación con este estudio & Limitación frente a esta SLR \\
\midrule
Priti y Kumar (2022) \cite{ART01} & Revisión crítica de modelos AQI (1960--2021). & Delimita el problema de comparabilidad entre índices y la heterogeneidad metodológica del campo. & No aborda inferencia difusa ni monitoreo en tiempo real como foco de análisis. \\
Suman (2020) \cite{ART02} & Revisión de métodos para interpretar el estado de la calidad del aire mediante AQI. & Refuerza la función comunicativa del AQI y su relación con la interpretación del riesgo. & No analiza decisiones de diseño difuso ni arquitecturas de monitoreo inteligente. \\
Murena (2004) \cite{PROB08} & Desarrollo y aplicación de un índice de contaminación en contexto urbano. & Acerca la discusión del AQI a un caso de implementación sobre área urbana amplia. & No incorpora inferencia difusa ni una lectura comparativa de plataformas en tiempo real. \\
Seibert et al. (2024) \cite{ART04} & Clasificación difusa de calidad del aire. & Representa la línea de uso de reglas difusas para clasificar estados de calidad del aire. & Se concentra en una tarea específica de clasificación y no en una síntesis de configuraciones de diseño. \\
Andrianto et al. (2024) \cite{ART06} & Monitoreo IoT con lógica difusa Mamdani e interpretación tipo ISPU/AQI. & Aporta un antecedente cercano a la integración entre sensores, índice e interpretación en tiempo real. & Corresponde a una implementación puntual y no permite por sí sola comparar variantes metodológicas del campo. \\
Lingaraj et al. (2025) \cite{ART08} & Arquitectura WSN optimizada para monitoreo de aire en smart cities. & Introduce la dimensión de operación distribuida y eficiencia de red en despliegues urbanos. & Privilegia la arquitectura de red y no articula de forma transversal variables, diseño difuso y evaluación del sistema. \\
\bottomrule
\end{longtblr}

\end{adjustwidth}
La comparación resumida en la Tabla~\ref{tab:antecedentes_intro} muestra que la dificultad actual proviene de la forma fragmentada en que la evidencia ha quedado reportada entre revisiones de AQI y desarrollos aplicados de monitoreo difuso. En los antecedentes revisados, las decisiones sobre variables de entrada, diseño del sistema difuso, arquitectura de monitoreo en tiempo real y evaluación del desempeño aparecen distribuidas entre estudios con alcances distintos y escasa posibilidad de lectura transversal \cite{ART01,ART02,PROB08,ART04,ART05,ART06,ART07,ART08}. Esa dispersión dificulta establecer qué configuraciones se repiten, qué decisiones metodológicas son comparables y qué criterios podrían servir para interpretar con mayor consistencia la evidencia disponible \cite{ART01,ART02,ART04,ART05,ART06,ART07,ART08}.

En atención a esa necesidad, este estudio se propone analizar de manera sistemática cómo se utilizan los sistemas de inferencia basados en lógica difusa en plataformas de monitoreo en tiempo real para evaluar la calidad del aire urbano. La pregunta de investigación que orienta la revisión es la siguiente: \textbf{¿Cómo se utilizan los sistemas de inferencia basados en lógica difusa en plataformas de monitoreo en tiempo real para evaluar la calidad del aire urbano y los riesgos asociados?}

A partir de esta pregunta, el resto del artículo se organiza de la siguiente manera. La Sección 2 describe la metodología de la revisión sistemática, incluyendo el diseño PICo, la estrategia de búsqueda, los criterios de selección, la evaluación de calidad y el esquema de extracción y síntesis. La Sección 3 presenta los resultados organizados según las directrices analíticas de la revisión. La Sección 4 discute los hallazgos, sus implicaciones y los vacíos que permanecen abiertos. Finalmente, la Sección 5 expone las conclusiones y las líneas de trabajo futuro.
